\documentclass[11pt,fleqn]{book}
\usepackage[utf8]{inputenc}
\usepackage[T1]{fontenc}
\usepackage{amsmath, amssymb, amsthm}
\usepackage{tcolorbox}
\usepackage{listings}
\usepackage{xcolor}
\usepackage{geometry}
\usepackage{titlesec}
\usepackage{hyperref}
\usepackage{graphicx}
\usepackage{enumitem}

% Page Geometry
\geometry{a4paper, margin=1in}

% Colors based on brute.ai theme
\definecolor{accent-primary}{HTML}{6366F1}   % Indigo
\definecolor{accent-secondary}{HTML}{10B981} % Emerald
\definecolor{highlight-intuition}{HTML}{AEBAFE}
\definecolor{highlight-maths}{HTML}{EA9C9C}
\definecolor{highlight-example}{HTML}{A0EFD8}
\definecolor{highlight-implementation}{HTML}{F1E881}
\definecolor{bg-secondary}{HTML}{F8FAFC}
\definecolor{border-color}{HTML}{E2E8F0}

% Custom TColorBox for Highlighting (mimicking the span.highlight)
\newtcbox{\highlightbox}[1][accent-primary]{on line,
  arc=3pt,colback=#1!20,colframe=#1!50,
  before upper=\rule[-3pt]{0pt}{10pt},boxrule=0.5pt,
  boxsep=0pt,left=2pt,right=2pt,top=2pt,bottom=2pt}

% Callout Style
\newtcolorbox{callout}[1][accent-primary]{
  colback=bg-secondary,
  colframe=border-color,
  leftrule=4pt,
  sharp corners,
  boxrule=0.5pt,
  title=Note,
  coltitle=accent-primary,
  fonttitle=\bfseries,
  attach title to upper,
  after title={:\enskip},
}

% Heading Style Enhancement
\titleformat{\chapter}[display]
  {\normalfont\huge\bfseries\color{accent-primary}}{\chaptertitlename\ \thechapter}{20pt}{\Huge}

% Listings Style for Python
\lstset{
  language=Python,
  backgroundcolor=\color{bg-secondary},
  basicstyle=\ttfamily\small,
  keywordstyle=\color{accent-primary},
  commentstyle=\color{accent-secondary},
  stringstyle=\color{orange},
  breaklines=true,
  frame=single,
  rulecolor=\color{border-color},
  padding=5pt
}

\title{Machine Learning Journal}
\author{brute.ai}
\date{March 2024}

\begin{document}

\maketitle
\tableofcontents

\chapter{Linear Regression Foundations}

Let's start our machine learning journey from this. This is the stepping stone or the early ideas of Machine learning. Simple as it sounds and ancestor of mathematical foundations in ML.

\begin{center}
    \textit{\textbf{In Linear Regression, a dependent variable $y$ is associated with set of independent variables $(x_1, x_2, ..., x_d)$ via a linear predictor function.}}
\end{center}

\textbf{For example}, let's take a popular housing price problem. So, we have a bunch of data points of house prices. Say, we know different properties of houses (\textit{In ML term, we call them features or attributes}). The features are area, number of bedrooms, kitchen size, garden, is parking available, location of house and so on. And we know the prices as well. So, typically the problem is given so many data points of house features and price, can we create a predictor function which can predict price for a new house given all the necessary features.

If we try to fit a linear function through the features, it will be a problem of \textbf{Linear Regression}.

\section{Intuition: Foundation of Linear Regression}

Let's define each term carefully. We assume to start with the housing price problem only. So, to make each prediction say we know $d$ different things about the house. Or in other words, our \textbf{number of features are $d$}.

\[X = \begin{bmatrix}
x_0 \\
x_1 \\
\vdots \\
x_{d-1}
\end{bmatrix}_{d \times 1}\]

So, our each independent variable matrix is like $X$. To compare with our example problem, $x_0$, $x_1$ are here bedroom size, area and so on. Now, we are trying to fit a linear function through the features. So, our predictor function is,
\[y = f_w(x) + \epsilon\]

Here $f_w(x) = w_0 + w_1x_1 + w_2x_2 + ... + w_dx_d$. $w_i$ is the weight associated with $i$-th feature (Also known as \textit{parameter}). $\epsilon$ is the irreducible error term. So, in matrix form $W$ is
\[W = \begin{bmatrix}
w_0 \\
w_1 \\
\vdots \\
w_{d-1}
\end{bmatrix}_{d \times 1}\]

\section{Pricing Example}
In the following example, we are trying to predict Rikshan price per distance travelled. The $x$ values are distance travelled (only one feature) and dependent predicted variables are $y$. We plot the data points and try to predict the ideal linear function to fit that data. Now for example, we assume the ideal function is \textbf{$y = w_0 + w_1x$}.

\begin{callout}
    To find the mathematically most accurate one, we need an \textbf{objective function} to reduce the cost of fit. For a complete new data point of distance we need to travel, we can find the $y$ value corresponding to that $x$ and assume that's closest possible prediction.
\end{callout}

\section{Irreducible Error}
If two variables are probabilistically related, then for a fixed value $x$, there is uncertainty in the value of the second variable. We assume the uncertainty as:
\[y = w_0 + w_1x + \epsilon\]
\textbf{$\epsilon$ is a random variable representing irreducible error}. Two variables $(y, x)$ are related linearly 'on average' if for fixed $x$ the actual value of $y$ differs from its expected value by a random amount (\textbf{random error}).
\[E(\epsilon) = 0, \quad V(\epsilon) = \sigma^2\]

\begin{callout}
    In Linear Regression, we assume that the irreducible noise comes from a \textit{Gaussian Distribution}. If we think of an entire population of $(x, y)$ pairs then $\mu_{y/x^*}$ is the mean of all $y$ values for which $x=x^*$ and $\sigma_{y/x^*}$ is a measure of how much these values of $y$ spread out about the mean value.
\end{callout}

\section{Estimating Parameters: The Closed Form Solution}
To find out the most suitable fit, we need a cost function to optimise. The most common cost function is \textbf{Least Squares Estimate}.
\[L = \sum_{(x_i,y_i \in D)} (y_i - f_w(x_i))^2\]

For a 1D case ($y = w_0 + w_1x$), the optimal parameters $w_0^*, w_1^*$ are found by solving:
\[w_1^* = \frac{\frac{\sum_{i}x_iy_i}{n} - \bar{X}\bar{Y}}{\frac{\sum_{i}x_i^2}{n} - \bar{X}^2}\]
\[w_0^* = \bar{Y} - w_1^*\bar{X}\]

In matrix form, the solution is elegantly represented as:
\begin{center}
    \Large{\textbf{\color{accent-primary} $w^* = (X^TX)^{-1}X^TY$}}
\end{center}

\chapter{Optimization with Gradient Descent}

In the previous section, we have seen that we can write the linear regression problem in the form of minimising the cost function. So, the problem reduces to finding the minimum of the cost function. 

But why do it iteratively? Well, in the case of linear regression, we have a closed-form solution. But that is not always the case for more complex models.

\section{Intuition: The Hill Descent}
The idea of this algorithm is to leverage the \textbf{gradient} of the loss function to find the optima. Consider the MSE loss function, which is convex and shaped like a bowl (parabola).

The gradient represents the slope of the curvature at a given point. If the gradient is positive, it means the function is increasing; if negative, it is decreasing. \textbf{Gradient Descent} tells us to take a step \textit{opposite} to the gradient to move towards the minima.

\section{The Update Rule}
For each dimension $j$, weights are updated iteratively:
\[w^{(j)} = w^{(j)} - \eta \frac{\partial L}{\partial w^{(j)}}\]
Where $\eta$ is the \textbf{learning rate} (a hyperparameter). For linear regression, this becomes:
\[w^{(j)} = w^{(j)} - \eta (E(w,x) - y) \cdot x^{(j)}\]

\section{Variants of Gradient Descent}
\begin{itemize}
    \item \textbf{Batch Gradient Descent}: Uses the entire dataset to compute gradients in every step. Guaranteed to converge for convex functions but slow for large data.
    \item \textbf{Stochastic Gradient Descent (SGD)}: Updates weights for every single training training example. Much faster and memory-efficient.
    \item \textbf{Mini-Batch Gradient Descent}: A compromise that updates weights based on a small subset (Batch Size $B$) of parameters.
\end{itemize}

\section{Advanced Optimizers}
\begin{itemize}
    \item \textbf{RMSProp}: Uses a moving average of squared gradients to normalize gradients, helping with vanishing/exploding gradients.
    \item \textbf{Adam}: Combines RMSProp with momentum (averaging the gradients themselves) for faster and more stable convergence.
\end{itemize}

\section{Implementation in Code}
\begin{lstlisting}[language=Python]
import torch
import torch.nn as nn

# Same thing in 3 lines using PyTorch
model     = nn.Linear(1, 1)
optimizer = torch.optim.SGD(model.parameters(), lr=0.1)
criterion = nn.MSELoss()

for epoch in range(500):
    optimizer.zero_grad()
    loss = criterion(model(X_t), y_t)
    loss.backward()          # computes all gradients automatically
    optimizer.step()         # w = w - lr * grad
\end{lstlisting}

\section{Probabilistic Interpretation of Least Squares}
Let us assume the target variables and the inputs are related via:
\[y_i = w^Tx_i + \epsilon_i\]
where $\epsilon_i$ is an i.i.d. random variable representing noise from a \textbf{Gaussian distribution} $N(0, \sigma^2)$. The density of $y_i$ given $x_i$ and weights $w$ is:
\[p(y_i|x_i,w) = \frac{1}{\sqrt{2\pi\sigma^2}} \exp\left(-\frac{(y_i - w^Tx_i)^2}{2\sigma^2}\right)\]

The \textbf{Maximum Likelihood Estimation (MLE)} of $w$ is obtained by maximizing the log-likelihood $LL(w)$:
\[LL(w) = C_1 - C_2 \sum_{i=1}^{n}(y_i - w^Tx_i)^2\]
Maximizing this is mathematically equivalent to minimizing the sum of squared errors.

\section{Maximum Aposteriori Estimation (MAP)}
In the Bayesian approach, we treat parameters as random variables and impose a \textbf{prior distribution} $p(w)$. The posterior probability is:
\[p(w|D) \propto p(D|w)p(w)\]
If we assume a Gaussian prior $p(w) = N(0, \frac{1}{\lambda} I)$, the MAP estimate is:
\[W_{MAP} = \arg\min_{w} \left[ \frac{1}{2\sigma^2} \sum_{j}(y_i-W^Tx_i)^2 + \frac{\lambda}{2} \|w\|_2^2 \right]\]
This introduces an additional term $\|w\|_2^2$, which corresponds to \textbf{L2 Regularisation}.

\chapter{The Bias-Variance Tradeoff}

\section{Train, Test, and Validation Splits}
To ensure a model generalizes well, we divide the dataset $D$ into:
\begin{itemize}
    \item \textbf{Training Set}: Used for model learning.
    \item \textbf{Testing Set}: Unseen data used for unbiased performance evaluation.
    \item \textbf{Validation Set}: Used to tune hyperparameters like learning rate.
\end{itemize}

\section{The Concept of the Tradeoff}
As model complexity increases (e.g., higher polynomial degree), \textbf{Training Error} always decreases. However, \textbf{Testing Error} typically follows a U-shaped curve: it decreases initially as the model learns patterns, but starts increasing as the model begins to overfit to the training noise.

\section{Mathematical Breakdown of Error}
Expected test error for a point $\hat{x}$ can be decomposed as:
\begin{center}
    \Large{\textbf{\color{accent-primary} $Err(\hat{x}) = \text{Bias}^2 + \text{Variance} + \sigma^2$}}
\end{center}
\begin{enumerate}
    \item \textbf{Bias}: Error from approximating a complex real-world problem with a simpler model.
    \item \textbf{Variance}: Error from the model's sensitivity to small fluctuations in the training set.
    \item \textbf{Irreducible Error ($\sigma^2$)}: Noise inherent in the data itself.
\end{enumerate}

\chapter{Decision Trees}

Decision Trees are rule-based models that split the data into branches to refine predictions.

\section{Splitting Criteria: Entropy and Information Gain}
\textbf{Entropy} $H(Y)$ measures the uncertainty or "chaos" in a node's class distribution:
\[H(Y) = -\sum_{i=1}^{k} P(y_i) \log_2 P(y_i)\]

\textbf{Information Gain} is the reduction in entropy after a split:
\[I(X, Y) = H(Y) - H(Y|X)\]
The goal is to choose the feature that maximizes this gain.

\section{Alternative: Gini Impurity}
The \textbf{Gini Index} is a computationally simpler measure of impurity:
\[Gini(t) = 1 - \sum_{i=1}^{k} p_i^2\]

\section{Combatting Overfitting}
\begin{itemize}
    \item \textbf{Pruning}: Cutting back branches that add little predictive power.
    \item \textbf{Thresholds}: Limiting tree depth or minimum samples per leaf.
\end{itemize}

\chapter{Recurrent Neural Networks (RNN)}

RNNs are designed for sequential data where the order of inputs matters (e.g., text, time series).

\section{The Core Architecture}
At each time step $t$, the RNN maintains a \textbf{hidden state} $h_t$, which acts as the "memory" of previous tokens:
\[h_t = \sigma(W_{hh}h_{t-1} + W_{xh}x_t + b_h)\]
\[y_t = W_{hy}h_t + b_y\]

\section{Applications}
RNNs (along with variants like LSTMs and GRUs) formed the foundation for modern NLP tasks, providing a way for neural networks to handle variable-length sequences.

\begin{callout}
    While RNNs face challenges like vanishing gradients, they represent the first major step humanity took toward solving complex sequential problems.
\end{callout}

\chapter*{Editor's Note}
This journal serves as a foundational guide to the mathematical and intuitive structures of Machine Learning. For interactive simulations and code walkthroughs, visit \url{https://brute.ai}.

\end{document}
